%! Function Fundamentals: Notation, Domain, and Range — Unit 2, Lesson 2.1 — Lesson Plan
\documentclass[10pt]{article}
\usepackage{atda-article}
\usepackage{atda-boxes}
\usepackage{graphicx}   % -article does NOT load graphicx; needed for thumbnails/screenshots
\graphicspath{{images/}}
\newcommand{\TallMath}[1]{$\displaystyle #1\rule[-1.4em]{0pt}{3.2em}$}
\newcommand{\UnitNumberName}{Unit 2: Functions, Transformations, and Graph Analysis \quad}
\newcommand{\LessonNumberName}{Lesson 2.1: Function Fundamentals: Notation, Domain, and Range}

\begin{document}
\begin{center}
    {\Large\bfseries \CourseName \\
    {\small\bfseries \UnitNumberName \LessonNumberName}}
\end{center}

\begin{tcolorbox}[colback=mist, colframe=royal, boxrule=0.9pt, arc=2mm,
    left=2mm, right=2mm, top=1.5mm, bottom=1.5mm]
    \textbf{Primary Objective:} I can read a graph in both directions --- given an input, find the
    output $f(x)$; given an output, find the input that produced it --- and write either read as a
    point $(x, f(x))$. I can state a function's domain and range from its graph, with units, in
    inequality, set, and interval notation, and explain what cuts a domain short: the situation,
    the graph, or the algebra.
\end{tcolorbox}

\renewcommand{\arraystretch}{1.6}
\begin{skillbox}[Priority Ideas \& Skills]{goldbox}
\begin{tabularx}{\linewidth}{@{}X|X@{}}
\begin{minipage}[t]{\linewidth}\vspace{0pt}
\textbf{Standards covered:}
\begin{itemize}[leftmargin=1.1em, itemsep=2pt, topsep=2pt]
  \item \texttt{AFDA.AF.2a} --- determine the domain and range of a function given a graphical
        representation, \emph{including those limited by contexts}.
  \item \texttt{AFDA.AF.2e} --- describe and recognize the connection between points on the graph
        and the value of a function.
\end{itemize}
\textbf{Skills students leave with:}
\begin{itemize}[leftmargin=1.1em, itemsep=2pt, topsep=2pt]
  \item Evaluate from a graph (\emph{input given}) and solve $f(x)=k$ from a graph
        (\emph{output given}) --- and know they are different motions.
  \item Translate freely among \textbf{sentence}, \textbf{function notation}, and \textbf{point}.
  \item State a domain and range in \textbf{inequality}, \textbf{set}, and \textbf{interval}
        notation, with units, including $\emptyset$.
  \item Name which of three things restricted a domain: the context, the graph, or the algebra.
  \item Justify why two inputs may share an output but one input may not have two.
\end{itemize}
\end{minipage}
&
\begin{minipage}[t]{\linewidth}\vspace{0pt}
\textbf{Key Understandings (the \emph{why}):}
\begin{itemize}[leftmargin=1.1em, itemsep=2pt, topsep=2pt]
  \item \textbf{``$f(a)=b$'' and ``$(a,b)$ is on the graph'' are one sentence in two costumes.}
        That equivalence is the whole of \texttt{AF.2e}, and it is what makes a graph readable
        at all --- every point is a completed function statement.
  \item \textbf{The notation is not symmetric.} $C(3)=10$ says nothing about $C(10)$. Students who
        stalled on ``solve $f(x)=1$'' in the Lesson 2.0 diagnostic are stalling here.
  \item \textbf{A domain is a decision, not a calculation.} The formula accepts $h=9$ happily; the
        garage does not. \texttt{AF.2a}'s phrase ``\emph{including those limited by contexts}'' is
        the standard telling you to teach exactly this.
  \item \textbf{A set of numbers has three standard spellings.} The AFDA guidance lists them
        explicitly --- equation/inequality, set notation, interval notation, and $\emptyset$ ---
        so all three are fair game on the test.
  \item \textbf{A function promises one output per input, and nothing in reverse.} The drone
        passes 24 feet twice. That is allowed, and saying why is the reasoning item of the day.
\end{itemize}
\end{minipage}
\end{tabularx}
\end{skillbox}

\begin{skillbox}[{Vocabulary, Concepts, \& Theorems}]{greenbox}
\renewcommand{\arraystretch}{1.35}
\begin{tabularx}{\linewidth}{@{}p{3.6cm}X@{}}
\textbf{Term} & \textbf{Meaning students record} \\
\hline
Function notation & $f(x)$ is the output the rule $f$ gives the input $x$. $f(a)=b$ and the point
    $(a,b)$ carry the same information. \\
Domain & Every input the function is allowed to take --- read left to right on the horizontal
    axis. \\
Range & Every output the function produces --- read bottom to top on the vertical axis. \\
Inequality notation & $0 \le h \le 8$. The everyday spelling; usually the easiest to read off a
    graph. \\
Set notation & $\{h \mid 0 \le h \le 8\}$, read ``the set of all $h$ such that\dots'' Also the
    only way to write a list, $\{2,5\}$, or $\{\;\}$ for nothing. \\
Interval notation & $[0,8]$ --- \textbf{smaller number first}. A bracket includes the endpoint, a
    parenthesis excludes it, and $\infty$ never takes a bracket. \\
Empty set $\emptyset$ & The set with nothing in it. Used when no input (or no output) works. \\
Restricted domain & A domain cut short by the context, by where the graph stops, or by the
    algebra (a zero denominator --- Lesson 1.7's excluded value). \\
Discrete domain & A domain that is a list of whole numbers. Its graph is separate dots, and it
    \emph{cannot} be written in interval notation. \\
\end{tabularx}
\end{skillbox}

\begin{fixedskillbox}[Activate Prior Knowledge \& Spiral Review]{mist}
\begin{tabularx}{\linewidth}{@{}X c@{}}
\begin{minipage}[t]{0.55\linewidth}\vspace{0pt}
\textbf{The warm-up is five items, and items 1 and 4 are the lesson in miniature:}
\begin{enumerate}[leftmargin=1.4em, itemsep=1pt, topsep=2pt]
  \item $f(x)=4x-9$: evaluate $f(5)$, then \emph{solve} $f(x)=7$ --- both directions
  \item An excluded value of $\frac{2x}{x+6}$ (Lesson 1.7; feeds notes box 4)
  \item Writing ``from $-2$ up to but not including $5$'' as an inequality (feeds notes box 3)
  \item The same both-directions question read off a \emph{table} --- and $g(x)=5$ has two answers
  \item Evaluate and interpret in context: a gym's $T(m)=15m+25$
\end{enumerate}
\end{minipage}
&
\begin{minipage}[t]{0.38\linewidth}\vspace{0pt}
\textbf{How to use the results:} item 1's second half and item 4's second half are the same skill.
A student who gets the first but not the second belongs in the Tier R group. Item 4 also plants
``two inputs, one output'' before the drone does --- do not resolve it during the warm-up, just
note who spotted it.
\end{minipage}
\end{tabularx}
\end{fixedskillbox}

\begin{skillbox}[Hook]{mist}
\textbf{The downtown garage.} \$4 to enter, \$2 an hour billed by the minute, and you may not stay
past 8 hours. One graph is posted by the gate.
\textbf{Pose it this way:} ``Two people walk up to the same sign. One asks \emph{what will three
hours cost me}. The other asks \emph{I only have \$16 --- how long can I stay}. Same graph. Are
you reading it the same way both times?'' Take answers before revealing the graph; most students
will say yes. Then reveal it and trace both reads with a finger --- up-then-across for the first,
across-then-down for the second --- and let the class name the difference themselves. Close the
hook by writing $C(3)=10$ and $(3,10)$ on the board side by side and saying they are the same
sentence. That equivalence \emph{is} \texttt{AFDA.AF.2e}, and everything else today is bookkeeping
on top of it.
\end{skillbox}

\boxguard
\begin{skillbox}[Lesson]{mist}
\begin{multicols}{2}
\textbf{1. Fill the notation table off the graph (10 min).} Four rows, and work them aloud in this
order: words $\to$ notation $\to$ point. Row 2 is the one that matters --- ``I have \$16'' becomes
$C(6)=16$, not $C(16)=6$. \textbf{Ask:} ``Which number went inside the parentheses, and how did you
decide?''

\textbf{2. Name the two motions (4 min).} Given the input, start on the horizontal axis and go up.
Given the output, start on the vertical axis and go across. Have students trace both on their own
copy. Then state the warning outright: $C(3)=10$ does \emph{not} mean $C(10)=3$.

\textbf{3. Domain and range off the same picture (6 min).} How far left to right, how far bottom
to top. \textbf{Ask:} ``Why doesn't the range start at \$0?'' (The entry fee is charged before the
clock starts --- the range's low end is a fact about the \emph{story}.) Then: ``Why does the domain
stop at 8 hours, when $2(9)+4=22$ is a perfectly good number?'' That question is
\texttt{AFDA.AF.2a} in one line.

\textbf{4. Three spellings of one set (10 min).} Work the notation table. The three rules that
settle everything: bracket includes, parenthesis excludes, $\infty$ never gets a bracket. Then the
fourth: an interval is written \emph{smaller number first}, no matter which way the graph runs ---
a candle burning down still has range $[0,20]$. Finish with the two rows that are \emph{not}
intervals: the list $\{2,5\}$ and the empty set.

\textbf{5. What cuts a domain short (6 min).} Three reasons: the context, the graph, the algebra.
Land the third on Lesson 1.7's $\frac{x+1}{x-3}$ --- the excluded value the class already found is
now a hole in a domain. Add the whole-number case (a field trip cannot take $17.4$ students) and
say plainly that such a domain \emph{cannot} be written in interval notation.

\textbf{6. Guided practice --- the drone (8 min).} $A(2)=24$ read forward; then $A(t)=24$ read
backward gives \textbf{two} answers. \textbf{Ask:} ``Is that allowed?'' It is --- and the reason is
the definition: one output per \emph{input}. Close by checking the peak against the formula, so
students see the graph and the algebra agree.
\end{multicols}
\end{skillbox}

\boxguard
\begin{skillbox}[Explicit Instruction: The Two Reads --- and the Three Spellings]{mist}
\begin{tabularx}{\linewidth}{@{}X|X@{}}
\begin{minipage}[t]{\linewidth}\vspace{0pt}
\textbf{The routine students record:}
\begin{enumerate}[leftmargin=1.4em, itemsep=2pt, topsep=2pt]
  \item \textbf{Which was I handed --- an input or an output?} That single question chooses the
        motion.
  \item \textbf{Input given:} horizontal axis $\to$ up to the graph $\to$ across. Write $f(a)=b$.
  \item \textbf{Output given:} vertical axis $\to$ across to the graph $\to$ down. Write $f(a)=b$
        again --- same statement, found from the other end.
  \item \textbf{Write the point} $(a,b)$ every single time. It is free, and it is the habit
        \texttt{AF.2e} is asking for.
  \item \textbf{For domain and range:} how far left to right, how far bottom to top --- then say
        it in all three notations, with units.
\end{enumerate}
\end{minipage}
&
\begin{minipage}[t]{\linewidth}\vspace{0pt}
\textbf{The three errors to name before they happen:}
\begin{itemize}[leftmargin=1.1em, itemsep=2pt, topsep=2pt]
  \item \textbf{Flipping the notation.} ``$L(6)=8$, so $L(8)=6$.'' Fix by pointing at the
        parentheses and saying \emph{the input lives here} every time. This is Tier E item 1(a).
  \item \textbf{Writing an interval backwards.} A decreasing graph tempts $[20,0]$. Interval
        notation is always low-to-high, regardless of direction. Tier E item 1(b).
  \item \textbf{Bracketing infinity.} $(3,\infty]$ appears every year. $\infty$ is a direction,
        not a number, so there is nothing to include.
\end{itemize}
\textbf{Say this out loud:} ``A function promises each input one output. It promises \emph{nothing}
about how many inputs share an output.''
\end{minipage}
\end{tabularx}
\end{skillbox}

\begin{skillbox}[Active Monitoring]{redbox}
\begin{itemize}[leftmargin=1.2em, itemsep=2pt, topsep=2pt]
  \item \textbf{Warm-up, items 1 and 4:} a student who evaluates fine and stalls on ``solve
        $f(x)=7$'' has the exact gap this lesson closes. Note the names before the notes start.
  \item \textbf{Notes, box 1, row 2:} listen for ``$C(16)=6$.'' Do not correct it with a rule ---
        ask ``is 16 dollars or hours?'' and let the units decide it.
  \item \textbf{Notes, box 3:} watch for brackets on $\infty$ and for intervals written high-to-low.
        Both are cheap to catch here and expensive on the test.
  \item \textbf{Activity, Tier A item 4:} groups often say ``there is no such height.'' Push them
        to the bottom of the bowl --- $H=0$ at $x=6$ and nowhere else.
  \item \textbf{Cold-call prompts:} ``Was I handed an input or an output?'' ``Which axis does the
        domain live on?'' ``Name the three things that can cut a domain short.'' ``Two inputs,
        one output --- allowed or not?''
\end{itemize}
\end{skillbox}

\boxguard
\begin{skillbox}[Group Work \& Differentiation]{redbox}
\begin{multicols}{3}
\textbf{Tier R --- Remediate}
\begin{itemize}[leftmargin=1em, itemsep=1pt, topsep=2pt]
  \item A burning candle, $L(t)=20-2t$: read $L(6)$, then solve $L(t)=4$ and check against the
        formula.
  \item Domain and range in two notations, with units.
  \item Why the domain stops at $t=10$ although $L(12)=-4$ computes fine.
\end{itemize}

\columnbreak
\textbf{Tier A --- Approaching Proficiency}
\begin{itemize}[leftmargin=1em, itemsep=1pt, topsep=2pt]
  \item A skate bowl, symmetric: read $H(3)$, then find \emph{both} $x$ with $H(x)=3$.
  \item Domain and range in two notations.
  \item The one height reached by a single input ($H=0$ at $x=6$), and what $H(15)$ means.
\end{itemize}

\columnbreak
\textbf{Tier E --- Extension}
\begin{itemize}[leftmargin=1em, itemsep=1pt, topsep=2pt]
  \item Critique two errors: flipped notation, and an interval written backwards.
  \item A \textbf{dot} graph (hot dogs): domain and range in set notation, and why an interval is
        the wrong tool.
  \item Justify: is the bowl not a function because one height happens twice?
\end{itemize}
\end{multicols}
\end{skillbox}

\begin{skillbox}[Individual Work \& Assessment]{redbox}
\textbf{Exit ticket (3 items, no notes):} a draining pool, $W(t)=2400-300t$ shown as a graph ---
(1) read $W(2)$ forward, as a point and as a sentence; (2) solve $W(t)=600$ and interpret it;
(3) \textbf{the conceptual check} --- domain and range in two notations, then \emph{why $t=10$ is
not in the domain even though $2400-300(10)$ is a number you can compute}.

\textbf{Using the results:} sort into three piles --- \emph{both directions clean},
\emph{forward only}, \emph{cannot start}. The middle pile is the largest and is the Tier R group
for Lesson 2.2. Item 3's second half separates students who think a domain is something you
calculate from students who understand it is something the situation decides; that distinction is
worth more than the notation.
\end{skillbox}

\boxguard[18]
\begin{skillbox}[Reinforcement \& Extension]{goldbox}
\textbf{Homework} (10 items): five fluency items (notation both ways, a table read both ways,
interval $\leftrightarrow$ inequality in each direction, and a three-part ``what restricts this
domain?''), then a four-part read of a highway arch --- the first graph students meet with
\textbf{negative inputs} --- ending in a clearance decision justified from the graph.

\textbf{Item 10 is the bridge to Lesson 2.2 and is worth grading closely.} Students tabulate
$f(x)=x^2$ against $g(x)=x^2+3$, see every output rise by $3$, and discover that the \emph{range}
moved from $[0,\infty)$ to $[3,\infty)$ while the domain did not move at all. They describe the
shift in their own words; \emph{translation} is 2.2's word to introduce, not today's to demand.

\textbf{Extension (optional):} Lesson 1.7's booster-club average cost $A(x)=\frac{6x+240}{x}$,
now asked as a domain question. Its domain is restricted \emph{twice} --- the algebra forbids
$x=0$ and the context forbids fractions and caps the order at 200 --- so the answer is
$\{1,2,3,\dots,200\}$ and interval notation cannot express it.

\textbf{Preview of Lesson 2.2 (\texttt{AFDA.AF.1a--b}):} parent functions and transformations ---
naming the move that item 10 just described.
\end{skillbox}

% ── Teacher notes — one per component, in packet order ────────────────────────
% Teacher-only prose lives HERE, never in a _key: a note is the one block in a
% key with no counterpart in the blank, so it makes the key longer than the
% blank. See references/conventions.md ("teachernote").
\boxguard[18]
\begin{teachernote}[Warm-Up]
    \textbf{8 minutes, then score it together.} Answers: (1) $f(5)=11$, $x=4$; (2) $x=-6$;
    (3) $-2 \le x < 5$; (4) $g(-1)=1$, and $g(x)=5$ at \emph{both} $x=-2$ and $x=2$;
    (5) $T(6)=115$, i.e.\ \$115 paid in all after six months. Items 1 and 4 are deliberately the
    same question in two representations --- a student fluent on the formula and stuck on the table
    is reading the table one column at a time instead of scanning the output row. Item 4's second
    half is the first appearance of ``two inputs, one output''; \textbf{do not resolve it during
    the warm-up.} Let it sit until the drone in guided practice, then point back. Item 2 is
    Lesson 1.7 returning as a domain question in notes box 4.
\end{teachernote}

\boxguard[18]
\begin{teachernote}[Guided Notes]
    \textbf{Box 1 is the lesson; protect its time.} Do not let the vocabulary box become a copying
    exercise --- fill \emph{function notation} and \emph{domain} together, assign the rest to pairs.
    In box 1, work row 2 out loud: ``I have \$16'' is an \emph{output}, so 16 goes after the equals
    sign. Insist on the point in the third column every time, even when it feels redundant; that
    habit is \texttt{AF.2e}. Box 3 is the SOL-notation box and it is dense --- work rows 1 and 2
    together and assign rows 3--5, then read the four rules aloud as a class. The two rows that are
    \emph{not} intervals ($\{2,5\}$ and the empty set) are the ones students skip, so ask for them
    by name. Box 4's three reasons are worth memorizing as a list; the whole-number paragraph
    underneath is what Tier E item 2 tests. In guided practice, item 2 is the reasoning item ---
    do not let the class answer ``because it's a parabola.'' Push for: at one instant the drone is
    at one altitude, but one altitude can happen at two instants. If time runs short, assign box 4
    as reading, but never cut box 1.
\end{teachernote}

\boxguard[18]
\begin{teachernote}[Group Activity]
    \textbf{Groups of three, 15 minutes, all groups start at Tier R.} Tier R is short by design; a
    group that clears it in four minutes should move straight on. Expect Tier R item 3 to produce
    inequality notation only --- ask for the interval as well, every time. Tier A item 2 is where
    groups slow down: many read only $x=3$ and stop, because they are scanning left to right and
    quit at the first hit. Prompt with ``is the bowl symmetric?'' Tier A item 4 (the single input
    at the bottom, $H=0$ at $x=6$) is the subtle one and is worth a share-out. Tier E item 1 is the
    share-out that matters --- read both wrong answers aloud and have the class diagnose them
    before the group explains. Tier E item 2(b) wants a \emph{specific} number that the interval
    allows and the situation does not: accept $2.5$ hot dogs or a \$7 price, but do not accept
    ``because you can't have decimals'' without a number attached. Tier E item 3 is the definition
    stated properly; a group that gets it can be asked to write the sentence on the board.
\end{teachernote}

\boxguard[18]
\begin{teachernote}[Exit Ticket]
    \textbf{5 minutes, notes closed, hand it to me at the door.} Answers: (1) $W(2)=1800$, the
    point $(2,1800)$ --- 1800 gallons left two hours in; (2) $t=6$ hours, with 600 gallons still
    to pump; (3) domain $0 \le t \le 8$ hours $=[0,8]$, range $0 \le W \le 2400$ gallons
    $=[0,2400]$, and $t=10$ is out because the pool is already empty at $t=8$ --- the formula's
    $-600$ gallons is not a thing that exists. Item 3's second half is the one to read carefully:
    a student who writes ``because the graph stops there'' has the graph reason, which is
    acceptable, but push in Lesson 2.2 for the \emph{context} reason underneath it. Do not grade
    this for points; use it to build the 2.2 groups.
\end{teachernote}

\boxguard[18]
\begin{teachernote}[Homework]
    \textbf{Expect 25--30 minutes.} Items 3 and 4 are the same translation in opposite directions
    and should be graded as a pair --- a student who can do 3 but not 4 has memorized a symbol
    rather than a meaning. Item 5 is the one to check first when planning the next do-now: 5(b)
    requires deciding that a tile's side is positive, which is a judgement, not a computation, and
    it is where most of the class will lose a point. The arch (6--9) is the first graph in this
    course with \textbf{negative inputs}; expect some students to report the domain as
    $0 \le x \le 10$ because they read only the right half. Item 9 must be justified from a value
    on the graph --- ``it looks tall enough'' earns nothing. Item 10 is the Lesson 2.2 bridge
    described in the plan above; accept any plain-English description of the upward shift. The
    extension is genuinely optional --- assign it to groups that finished Tier E early, and note
    that its domain is restricted \emph{twice}, which is the first time students have met that.
\end{teachernote}
\end{document}
